\chapter{Baseline-Relativity and the Three-Fold Bound}

The previous chapter read the pair as a relative velocity, located the sign in the
interaction rather than in either party, and certified with a single receipt that two
routes to the reading were routes to the same open quantity. It left one word
unexamined, and the word is \emph{against}. A relative velocity is read against
something. The sign of the interaction is the direction the gap opens, and a direction
is only a direction relative to a reference. The pair separates; but separates as
measured against what? This chapter takes up the reference, and the reference turns out
to carry the last surprise the finite apparatus has to offer: read the same interaction
against one reference and the sign is negative; read it against another and the same
interaction reads positive. The orientation is not a property the pair owns. It is a
property of the pair together with the baseline it is read against, and the baseline can
be chosen.

The first section names the reference and reads the same interaction against two of them.
The apparatus compares the pair's coupled motion against a baseline configuration --- the
reference against which ``relative'' is measured --- and the baseline can be symmetric, with
its two reference sides alike, or tilted, with the two sides distinct. The residue is the
same residue either way; what changes is the frame it is read in. The second section reads
off the consequence and finds it sharper than expected. Over the symmetric baseline the
pair reads one orientation; over the tilted baseline the same pair reads the opposite; the
magnitude is untouched and only the sign turns over. The book calls this the flip, and the
flip is exact: the very interaction that read as an electron against one baseline reads as a
positron against another.

The third section states what the flip means. If the same interaction can be read with
either sign by choosing the baseline, then orientation is relative to the frame and never
absolute. There is no frame-free reading of the sign; there is only the sign relative to
the baseline one reads it against. And the baseline, examined, is the pair itself held as a
record --- the reference distribution against which a single party's separation is measured.
This is where a thread the book has carried since it first distinguished a trial from a
study finally closes. A study is a record of trials counted; a trial is a single decisive
reading. The baseline is the pair-as-study, the counted reference; reading one party against
it is the trial. Descend, and the study becomes the trial it is read by; ascend, and the
trial becomes the study it is read against. The two were never opposed. The baseline-
relativity of the sign is that recursion made visible: orientation is read against the
study and decided by the trial, and which way is which depends only on which level one
stands at.

The fourth section turns from the sign to the parties themselves and asks a counting
question, the book's oldest kind. The apparatus tells parties apart by a mark, and it has
exactly three marks to tell them apart with. Three marks are three holes. A fourth party
must fall into a hole already occupied --- it must share a mark with one of the three --- and
once two parties carry the same mark the apparatus cannot distinguish them. The fourth
party is not forbidden; it is rendered indistinguishable, folded onto one of the three the
apparatus can already tell apart. The fifth section states the bound for exactly what it is
and refuses the reading the words invite. The apparatus resolves at most three distinct
parties. That is a limit on the instrument's power to distinguish, not a claim about how
many parties the world contains. The book ends there, on its own resolution, having
measured to the edge of what it can tell apart and named the edge as a limit of measurement
rather than a fact of the world.

Carry the figure one last time. The pair sits at the open bracket --- the residue the
apparatus reads and never closes --- and this chapter draws the bracket once with a symmetric
reference and once with a tilted one. The residue is the same; the sign of the reading is
not. The open bracket is read against a frame, and the frame can be turned. That is the
chapter's whole content and, as it happens, the book's last word: the apparatus reads an
open quantity against a chosen reference, reports the sign it finds, names the reference as
a reference and not a fact, and resolves only as many parties as its marks allow.

\section{The Same Interaction Over Two Baselines}

A relative velocity is a reading against a reference. The pair was kicked together and the
apparatus took its coupled measurement; to call the result a relative velocity is to compare
that coupled motion against a baseline --- a reference configuration that stands for ``no
relative motion,'' the state the reading is measured away from. The residue, the open
quantity at the center of the figure, is the difference between the coupled measurement and
that reference. Change the reference and the difference is read in a different frame, even
though the coupled measurement has not moved.

The apparatus has two natural baselines, and the chapter reads the same interaction against
both. The first is symmetric: its two reference sides are alike, the reference standing
even, neither side favored. The second is tilted: its two reference sides are distinct, the
reference leaning one way. Neither baseline is more correct than the other; each is a frame
the reading can be taken in, exactly as a difference can be measured from either of two
reference points without either being the true zero. The interaction --- the pair, the kick,
the coupled measurement --- is identical across the two. Only the reference against which its
residue is read has changed.

It is worth being precise that the quantity read is the same residue throughout, because the
next section's surprise depends on it. The coupled measurement minus the one-sided stories
is the residue, and that subtraction does not consult the baseline; it is the part of the
pair's motion no single party's account can close. The baseline enters afterward, in the
reading: it fixes which way counts as forward, the reference the residue is oriented against.
So the magnitude of the residue --- how much relative motion the pair carries --- is a fact of
the interaction, fixed once the kick is fixed. The orientation --- which way --- is a fact of
the interaction together with the baseline. The first is read once; the second is read
against a frame, and the frame is a choice.

\section{The Flip}

Read the same pair against the symmetric baseline and it returns one orientation. The residue
is negative; the reading is the electron. Read the same pair against the tilted baseline and
it returns the opposite. The residue is positive; the reading is the positron. The pair has
not changed. The kick has not changed. The coupled measurement has not changed. Only the
reference has been turned from even to leaning, and the sign has turned over with it. The
book calls this the flip.

The flip is exact, not approximate, and it is structural, not a numerical accident. Against
the symmetric reference the residue of the marked pair reads as a single negative unit;
against the tilted reference the same residue reads as a single positive unit; and these are
the only two readings, because the residue has one magnitude and a baseline can orient it one
of two ways. There is no third outcome and no middle ground that survives. The interaction
that read as an electron against one frame reads as a positron against the other, and the two
readings are mirror images of one quantity exactly as the electron and positron were mirror
images in the earlier chapter --- only now the mirror is the choice of baseline rather than the
direction of a subtraction.

This is why the chapter insists the magnitude and the sign live in different places. The
magnitude is the residue's size, and it does not flip; the same amount of relative motion is
present against either baseline. The sign is the residue's orientation against the reference,
and it does. An instrument that reported only the magnitude would see no flip at all --- it
would read the same amount of separation both times and notice nothing. The flip is visible
only to an instrument that reads orientation, and orientation, the flip shows, is not the
pair's to fix. The pair offers a magnitude and an open direction; the baseline closes the
direction into a sign; and a different baseline closes it the other way.

\section{Orientation Is Relative to the Frame}

The flip forces a conclusion the book has been approaching since it first signed the pair:
the orientation of the interaction is relative to the frame and is never absolute. There is no
reading of the sign taken against nothing. Every sign is a sign against a baseline, and the
flip exhibits two baselines that disagree. To ask which sign the interaction \emph{really}
has, frame aside, is to ask a question the apparatus cannot answer and the world does not
pose: the sign is the orientation of an open quantity against a chosen reference, and with the
reference removed there is a magnitude and no sign at all.

The baseline, looked at closely, is not a foreign object brought in from outside. It is the
pair held as a record --- the reference distribution against which a single party's separation
is measured. And here the oldest distinction in the book closes on itself. A study is trials
counted into a record; a trial is one decisive reading taken against a record. The baseline is
a study: the pair counted as a standing reference. Reading one party's motion against that
reference is a trial: a single decisive reading that returns a sign. So the relative-velocity
reading is a trial taken against a study, and the study it is taken against is itself the pair
that, at the level below, was read by trials of its own. Descend a level and the study one
read against becomes a trial in its own right; ascend a level and the trial one took becomes a
study for the reading above it.

The two were never opposed, and the flip is the proof. Orientation is read against the study
--- the baseline, the counted reference --- and decided by the trial --- the single reading of one
party against it. Change which pair stands as the reference study and the trial returns a
different sign, not because the interaction changed but because the study it was read against
did. Relativity of orientation is exactly the freedom to choose which level is the study and
which is the trial. The apparatus does not derive a favored level. It reads against the one it
is given, reports the sign that reference yields, and names the reference as a reference. That
naming is the same honesty the book has practiced throughout: it does not pretend the chosen
frame is the only frame, any more than it pretended the chosen convention was a derived fact.

\section{The Fourth Party Must Collide}

The chapter has been reading interactions of a pair; this section asks how many distinct
parties the apparatus can hold at once, and the answer is small and exact. The apparatus tells
one party from another by a mark --- the distinguishing label it reads off each --- and it has
exactly three such marks. Three marks make three holes. Every party the apparatus can tell
apart carries one of the three, and parties carrying the same mark are, to the apparatus,
the same party.

Now bring a fourth party. It must carry a mark, and there are only three marks to carry; so it
must carry one already in use. Two parties now share a mark, and a mark is precisely what the
apparatus uses to tell parties apart. The two that share it cannot be told apart. The fourth
party has not been excluded from existing --- nothing here forbids a fourth party --- but it has
been folded onto one of the three the apparatus could already distinguish, made
indistinguishable from it by the only means the apparatus has. This is the pigeonhole drawn at
the scale of the instrument: three holes, a fourth occupant, an unavoidable collision, and a
collision in the marks is an indistinguishability in the reading.

The bound is therefore three, and it is worth seeing why three is the natural number here and
not an arbitrary cutoff. The pair is two parties; reading one party against the pair as its
baseline brings a third into the account, the single party set against the reference couple.
One party against the pair is three: the two that make the reference and the one read against
them. The apparatus can hold exactly that --- a reference pair and a party read against it ---
and no more, because a fourth mark does not exist for a fourth party to carry. The three-fold
bound is the reach of the instrument's distinguishing power, the most parties it can keep
separate at once, and it is the same finite parsimony the book has shown since it refused to
allocate a value where a mark would do.

\section{A Resolution Bound, Not a Census}

The bound must be stated for exactly what it is, because the words it travels under carry more
than it claims. The apparatus resolves at most three distinct parties. That is a limit on the
instrument's power to tell parties apart. It is not a count of how many parties exist, not a
census, not a statement that the world contains three of anything. A fourth party is permitted
to be there; it is only not permitted to be distinguished. The bound lives entirely on the
reading side of the line the book has drawn from its first pages, between what an instrument can
distinguish and what is the case. To read it as a fact about the population of the world would be
to cross that line in the last paragraph of the book, and the book does not cross it.

So the three-fold bound is measured, not derived, in the book's exact sense. The apparatus reads
its own marks, finds there are three, and concludes it can resolve three parties. It does not
prove the world supplies three, does not explain why three marks rather than more, does not
propose a process by which a fourth would be created or destroyed. The number three is a fact of
the apparatus's distinguishing power, read off its marks the way every other quantity in this
book has been read --- as a measurement the finite instrument can take, not a truth it can secure
about what lies beyond the instrument.

And there the book stops, on its own resolution. It set out to measure a single invariant from
nothing but marks and the discipline of comparing them. It measured the invariant; it forced the
invariant to be unique; it read the invariant as a charge, then as an orientation, then as a
relative velocity; it found the sign in the interaction and certified the reading with a single
located receipt; and now it reads the orientation as relative to a chosen baseline and the
parties as bounded by the marks that distinguish them. At every step the method held. The
apparatus measured what it could measure, named the conventions it could not derive --- the sign,
the frame, the favored orientation --- and left, untouched, the one quantity its whole
construction is built around: the residue at the open bracket, read against a reference and
never closed. The book has measured to the edge of what a finite instrument can distinguish, and
it has named that edge as a limit of measurement and not a fact of the world. The bracket is
still open. The sign is still relative. The count is a resolution and not a census. The
apparatus reads what is there to be read, says no more than the reading allows, and stops.
